\documentclass[journal]{IEEEtran}
\usepackage{cite}
\usepackage{amsmath,amssymb,amsfonts}
\usepackage{algorithmic}
\usepackage{graphicx}
\usepackage{textcomp}
\usepackage{xcolor}
\usepackage{booktabs}

\begin{document}

\title{Architecting a Pure P2P Protocol for Million-Scale\\Live Video Streaming: A Taxonomy of Challenges\\and Decentralized Solutions}

\author{Author Name(s)%
\thanks{The authors are with the Department of Computer Science...}}

\maketitle

\begin{abstract}
The centralized model of live video streaming, reliant on Content Delivery Networks (CDNs), presents significant scalability costs and single points of failure. While peer-to-peer (P2P) networks offer a compelling decentralized alternative, existing protocols fail to address the unique constraints of low-latency live broadcasting at a scale of millions of concurrent users. This paper presents a comprehensive taxonomy of the fundamental challenges in designing a purely decentralized, tokenless live streaming protocol capable of this scale. We systematically identify and analyze five core problems: (1) the Live Chunk Propagation Problem, where naive distribution leads to exponential latency; (2) the Massive Churn Problem, where peer volatility disrupts data flow; (3) the Decentralized Coordination Problem, for discovering and synchronizing the live data window; (4) the Leeching Problem, enforcing fairness without central authority; and (5) the NAT Traversal Problem, a pragmatic barrier to pure decentralization. For each challenge, we propose and analyze solution classes that adhere to a strict decentralized ethos, including self-organizing multi-tree topologies, gossip-based repair protocols, Proof-of-Upload (PoU) receipts, and a Tit-for-Tat strategy based on verifiable, gossiped reputation. This work provides a foundational framework and defines the critical research trajectory for achieving robust, million-scale P2P live streaming.
\end{abstract}

\begin{IEEEkeywords}
Peer-to-Peer, Live Streaming, Decentralized Systems, Protocol Design, Scalability, Churn, NAT Traversal, Tit-for-Tat.
\end{IEEEkeywords}

\section{Introduction}
Global live video traffic is experiencing unprecedented growth, driven by sectors including gaming, social media, and teleconferencing. This growth is predominantly serviced by a centralized Content Delivery Network (CDN) model, which incurs immense economic costs and creates inherent bottlenecks and points of failure \cite{smith2022economics}. Peer-to-peer (P2P) architectures, famously exemplified by BitTorrent \cite{cohen2003incentives} for static file distribution, offer a paradigm shift. By leveraging the upload capacity of consumers themselves, a P2P system can theoretically scale organically with demand, reducing reliance on central infrastructure.

The transition from distributing static files to streaming live video, however, introduces a set of stringent constraints that existing P2P protocols are ill-equipped to handle. Live streaming is a real-time service with strict latency requirements; a delay exceeding several seconds is often unacceptable. Furthermore, the user base is highly volatile, with peers frequently joining and leaving the stream. A naive application of a BitTorrent-like swarm to live data results in catastrophic failure at scale, characterized by unbounded latency and network collapse under churn.

This paper argues that the goal of a purely decentralized, tokenless protocol for million-scale live streaming is achievable only by directly addressing a well-defined set of interconnected challenges. We move beyond high-level description to provide a detailed taxonomy of these problems and analyze solution classes that operate within a strictly decentralized constraint. Our contribution is a framework that delineates the problems of Propagation, Churn, Coordination, Fairness, and Connectivity, and evaluates proposed mechanisms such as structured topologies, verifiable reputation, and gossip protocols as potential solutions. We conclude that the path forward requires a departure from unstructured flooding towards a hybrid push-pull model within a dynamically managed overlay network.

\section{Background and Related Work}
Early P2P live streaming systems such as CoolStreaming \cite{zhang2005coolstreaming} and PPLive \cite{pplive2005} demonstrated feasibility for thousands of users by employing mesh-pull architectures and data-driven overlay construction. However, these systems often relied on central trackers and were not designed for the million-user scale or the strict decentralization targeted here.

Modern approaches often incorporate blockchain or token-based incentives to reward uploaders \cite{livepeer2023}. While effective for incentivization, these systems introduce complexity, volatility, and a form of financial centralization that contradicts the goal of a purely tokenless protocol. Our work focuses on the underlying data distribution problem, assuming that technical solutions for fairness can be achieved without cryptographic tokens.

WebRTC has emerged as a key enabler, providing native browser-based P2P connections and NAT traversal capabilities. Projects like WebTorrent \cite{webtorrent} extend the BitTorrent protocol to use WebRTC transports, but they focus primarily on static file sharing and lack the mechanisms for low-latency live propagation. Our work seeks to adapt these foundational technologies specifically for the live streaming context.

This paper synthesizes ideas from these domains---structured overlays from early systems, verifiable accounting from blockchain-inspired designs, and modern transport protocols---to form a coherent blueprint for a next-generation, decentralized live streaming protocol.

\section{A Taxonomy of Challenges and Decentralized Solutions}
The design of a million-user P2P live stream is fundamentally constrained by five critical challenges. Table \ref{tab:taxonomy} summarizes these challenges and the proposed solution classes.

\begin{table*}[t]
\centering
\caption{Taxonomy of Challenges and Decentralized Solution Classes for Million-Scale P2P Live Streaming}
\label{tab:taxonomy}
\begin{tabular}{lp{6.2cm}p{7.8cm}}
\toprule
\textbf{Challenge} & \textbf{Description} & \textbf{Proposed Solution Classes} \\ \midrule
\textbf{1. Live Chunk Propagation} & 
The time-sensitive dissemination of sequential video chunks from a source to all peers without central fan-out. Naive flooding leads to exponential latency growth. & 
\textbf{Structured Multi-Tree Topology:} Actively construct a low-diameter distribution tree (or forest) where high-capacity ``superpeers'' are promoted close to the source. \newline
\textbf{Hybrid Push-Pull:} The source and upper-layer peers push chunks to their immediate children, who then serve downstream peers. \\ \midrule

\textbf{2. Massive Churn} & 
The constant joining and leaving of peers disrupts the distribution tree, causing data loss and buffering for downstream peers. & 
\textbf{Redundant Parent Connections:} Each peer connects to multiple ($k$) parents for the same data stream. \newline
\textbf{Gossip-Based Repair:} A subsidiary gossip protocol allows orphans (peers whose parent failed) to quickly discover and connect to candidate parent nodes. \\ \midrule

\textbf{3. Decentralized Coordination} & 
Peers must discover the existence of the live data window and coordinate playback without central trackers. & 
\textbf{Gossiped Announcements:} The source gossips a signed message containing the latest chunk ID (sequence number). This announcement propagates with a limited TTL, defining the ``live window.'' \newline
\textbf{Live DHT:} An extension of the Kademlia DHT \cite{mazieres2002kademlia} optimized for rapid lookups of the most recent chunks. \\ \midrule

\textbf{4. Leeching \& Fairness} & 
Self-interested peers (free-riders) consume network resources without uploading content, causing system-wide performance degradation. & 
\textbf{Tit-for-Tat (TFT):} Enforce direct resource exchange, prioritizing peers that actively upload to their neighbors. \newline
\textbf{Verifiable Reputation:} Use cryptographically signed Proof-of-Upload (PoU) receipts gossiped locally to rank peer contribution. \\ \midrule

\textbf{5. NAT Traversal} & 
Firewalls and symmetric NATs prevent direct peer-to-peer connections, disrupting the overlay formation. & 
\textbf{QUIC Hole Punching:} Use STUN/TURN with UDP-based QUIC transport for native, high-probability direct connectivity. \newline
\textbf{Emergent Relay Selection:} Autonomously promote nodes with public IP addresses and high capacity to act as localized relays. \\ \bottomrule
\end{tabular}
\end{table*}

\subsection{The Live Chunk Propagation Problem}
The core tension lies in reconciling the inherent propagation latency of P2P networks with the strict temporal constraints of live media. An unstructured swarm, where chunks are pulled at random, results in a propagation delay that grows logarithmically or linearly with the number of peers \cite{kumar2019chunk}. For a million peers, this delay becomes minutes, not seconds.

\subsubsection{Proposed Solution: Managed Topology}
The solution requires actively constructing a distribution topology rather than relying on an emergent mesh. The protocol must dynamically build a multi-tree structure where high-bandwidth, stable peers are promoted close to the video source (Layer 1 and Layer 2), while low-bandwidth consumers and leaf nodes occupy deeper layers. Combined with a hybrid push-pull scheduling policy, live chunks are pushed down pre-established paths of the tree structure for speed, while missing chunks are pulled from neighboring peers via local mesh connections to guarantee reliability.

\subsection{The Massive Churn Problem}
At a scale of millions, thousands of users join and leave every minute. Single-parent tree topologies suffer catastrophic failure as subtrees get disconnected when a parent departs.

\subsubsection{Proposed Solution: Gossip-based Repair \& Multi-Parent Mesh}
To mitigate peer volatility, each node must maintain multiple concurrent upstream providers (parents) rather than a single parent. If one parent disconnects, the data flow is maintained through the remaining parents. Concurrently, a lightweight, background gossip protocol continuously provides a list of candidate backup nodes, allowing immediate connection replacement of failed parents without relying on a tracker. Missing symbols are quickly requested via a swarm repair mechanism from local mesh neighbors.

\subsection{The Decentralized Coordination Problem}
Without a central tracker or signaling server, peers must self-discover active streams and align their playback windows.

\subsubsection{Proposed Solution: DHT Stream Registration \& Head Tracking}
A Distributed Hash Table (DHT), such as Kademlia, is utilized to store stream IDs mapping to active peer identities. To avoid DHT bottlenecks, the DHT is used only for bootstrap discovery. Once connected, peers use HyParView-style gossip protocols to maintain local membership. To synchronize playback, nodes periodically gossip their current playback "head" (e.g., chunk index 15000), allowing nodes lagging behind to aggressively fetch newer chunks and self-align to the live edge.

\subsection{The Leeching Problem}
In a pure P2P network, self-interested peers may download stream chunks without uploading any, exhausting the network's collective capacity.

\subsubsection{Proposed Solution: Verifiable Reputation \& Tit-for-Tat}
The system enforces resource contribution through a reputation-driven Tit-for-Tat mechanism. Peers issue signed Proof-of-Upload (PoU) receipts when receiving chunks from an uploader. These receipts form a verifiable, non-repudiable contribution history. This local reputation is gossiped among neighboring nodes. Nodes prioritize upload slots for peers with high reputation scores, while free-riders are systematically demoted to deeper layers or disconnected.

\subsection{The NAT Traversal Problem}
A substantial fraction of internet users reside behind symmetric NATs or firewalls, preventing direct incoming and outgoing P2P connections.

\subsubsection{Proposed Solution: Direct Hole Punching \& Emergent Relays}
Peers utilize modern UDP-based transport protocols like QUIC or WebRTC, integrated with STUN/TURN for direct hole punching. For peers where direct connection is mathematically impossible (e.g., dual symmetric NAT), the overlay promotes highly capable, publicly accessible super-peers as emergent relay nodes. Relaying is treated as a high-reputation upload activity, which incentivizes open nodes to provide routing capacity, ensuring complete coverage for NAT-restricted leaf nodes.

\section{Conclusion and Future Work}
This paper has presented a taxonomy of the critical challenges facing the design of a pure P2P live streaming protocol for millions of users. We have argued that solutions exist within a decentralized framework, centered on a verifiable reputation system and a self-organizing managed topology. The most significant remaining obstacle is the efficient and decentralized mitigation of the NAT traversal problem. Future work will involve the simulation and evaluation of the proposed hybrid push-pull model and the gossip-based reputation system to quantify the achievable latency and robustness under churn at massive scale.

\begin{thebibliography}{1}

\bibitem{smith2022economics}
J.~Smith, ``The Economics of Video Delivery,'' {\em IEEE Internet Comput.}, vol.~26, no.~3, pp.~45--51, May 2022.

\bibitem{cohen2003incentives}
B.~Cohen, ``Incentives Build Robustness in BitTorrent,'' in {\em Workshop on Economics of Peer-to-Peer Systems}, vol.~6, 2003, pp.~68--72.

\bibitem{zhang2005coolstreaming}
X.~Zhang, J.~Liu, B.~Li, and T.-S.~P. Yum, ``CoolStreaming/DONet: A Data-driven Overlay Network for Efficient Live Media Streaming,'' in {\em Proc. IEEE INFOCOM}, 2005, vol.~3, pp.~2102--2111.

\bibitem{pplive2005}
``PPLive: Internet Television via Peer-to-Peer,'' [Online]. Available: http://www.pplive.com.

\bibitem{livepeer2023}
``LivePeer: The Decentralized Live Video Streaming Protocol,'' LivePeer, Inc., White Paper, 2023.

\bibitem{webtorrent}
``WebTorrent,'' [Online]. Available: https://webtorrent.io.

\bibitem{mazieres2002kademlia}
P.~Maymounkov and D.~Mazières, ``Kademlia: A Peer-to-Peer Information System Based on the XOR Metric,'' in {\em Proc. IPTPS}, 2002, pp.~53--65.

\bibitem{kumar2019chunk}
A.~Kumar et al., ``Chunk Propagation in Live P2P Streaming,'' {\em ACM Trans. Multimedia Comput. Commun. Appl.}, vol.~15, no.~2s, pp.~1--24, Jun. 2019.

\end{thebibliography}

\end{document}
